% EXTRACTED FOR PAPER TWO (the hardware case study) -- PASS 48.1
% source     : paper/paper.tex
% blob       : 13590dbfe8119022ae95ff7a374910e30146ab8c
% lines      : 1435-1450 (1-indexed, inclusive)
% section    : Section 7
% prose words: 139
%
% Related Work: 'Statistical-to-bias-floor transitions on hardware' -- the March 2026 report of O(M^-1/2) decay saturating at a hardware floor, and how it differs from ours.
%
% This is a VERBATIM COPY. The original is unmodified and still frozen.
% Regenerate with: experiments/pass48_extract_paper2.py

\textbf{Statistical-to-bias-floor transitions on hardware.} A March 2026
study of shadow tomography on an integrated photonic processor
\cite{hardwarehorizon2026} reports reconstruction error following
\(O(M^{-1/2})\) and then saturating at a hardware-determined floor, a
transition the authors term a Hardware Horizon. The structure is adjacent
to Section 5 and the two should be distinguished: their transition is
between statistical scaling and a bias floor at fixed system size, ours is
a change in the statistical exponent itself as a function of system size,
occurring before any hardware bias enters. The rolling-exponent diagnostic
is likewise not ours: Aasen and G\"arttner \cite{aasen2026limitations}
fit a power-law exponent to adaptive-tomography infidelity in sliding
budget windows and report it drifting from \(1/N\) toward
\(1/\sqrt{N}\), but their migration is driven by detector noise eroding
an adaptive advantage, while the one in Section 3.6 is a property of the
noiseless estimator's own variance.

